%% 1-11 目もりを よもう  第2部の さいごの 章（第11回で 説明を ふやした）
%% 第1巻のきまり：15 pt・行送り28 pt、分かち書き（文節の間に半角の空白）、漢字はすべてルビ、式は全角の＋－＝で書く
%% この章の部品（名前は \bkqj…）
\makeatletter
% まっすぐな 目もり  \bkqjscale{(x,y)}{1目もりの長さmm}{目もりの数}{数字を書く間（目もりいくつごと）}{はじめの数}{数字の増え方}
%   (x,y) は いちばん左の目もり。目もりは上のふちから下へ、数字は下に書く。
\newcommand\bkqjscale[6]{\begin{scope}[shift={#1}]
  \pgfmathsetmacro\bkqj@L{#2*#3}
  \fill[bkYellowBg, draw=bkYellow!80!black, line width=0.7pt, rounded corners=0.8mm] (-5,-8) rectangle (\bkqj@L+5,3);
  \foreach \k in {0,...,#3} {\pgfmathtruncatemacro\bkqj@lab{mod(\k,#4)==0}
    \ifnum\bkqj@lab=1 \draw[bkInk, line width=0.8pt] (\k*#2,3) -- (\k*#2,-2.2);
      \pgfmathtruncatemacro\bkqj@v{#5+\k/#4*#6}
      \node[font=\sffamily\bfseries\footnotesize, text=bkInk] at (\k*#2,-5) {\bkqj@v};
    \else\draw[bkInk!85, line width=0.5pt] (\k*#2,3) -- (\k*#2,-0.2);\fi}
\end{scope}}
% 目もりを さす 赤い しるし（上から）  \bkqjmark{(x,y)}   (x,y) は 目もりの 上のふち
\newcommand\bkqjmark[1]{\fill[bkRed] #1 ++(0,0.6) -- ++(-1.8,3.4) -- ++(3.6,0) -- cycle;}
% はんえんの メーター  \bkqjgauge{(x,y)}{半径}{いちばん大きい数}{数字を書く間}{目もりの数}{針}
\newcommand\bkqjgauge[6]{\begin{scope}[shift={#1}]
  \fill[bkGray!35, rounded corners=2mm] (-#2-3,-5) rectangle (#2+3,#2+3);
  \fill[white] (-#2,0) arc (180:0:#2) -- cycle;
  \draw[bkInk!70, line width=0.8pt] (-#2,0) arc (180:0:#2);
  \foreach \k in {0,...,#5} {\pgfmathsetmacro\bkqj@a{180-\k/#5*180}\pgfmathsetmacro\bkqj@v{#3/#5*\k}%
    \pgfmathtruncatemacro\bkqj@lab{mod(round(\bkqj@v),#4)==0}\pgfmathtruncatemacro\bkqj@vi{round(\bkqj@v)}
    \ifnum\bkqj@lab=1 \draw[bkInk, line width=0.8pt] (\bkqj@a:#2*0.96) -- (\bkqj@a:#2*0.8);
      \pgfmathsetmacro\bkqj@fs{#2*0.36}%
      \node[font=\sffamily\bfseries\fontsize{\bkqj@fs}{\bkqj@fs}\selectfont, text=bkInk, inner sep=0pt] at ($(\bkqj@a:#2*0.64)+(0,1.2)$) {\bkqj@vi};
    \else\draw[bkInk!80, line width=0.4pt] (\bkqj@a:#2*0.96) -- (\bkqj@a:#2*0.87);\fi}
  \draw[bkRed, line width=1.2pt, line cap=round] (0,0) -- ({180-#6/#3*180}:#2*0.86);
  \fill[bkInk!80] (0,0) circle (1.3);
\end{scope}}
% 目もりの ある コップ  \bkqjcup{(x,y)}{はば}{1目もりの高さ}{目もりの数}{水の高さ（目もりで）}   (x,y) は底のまん中。左に数字
\newcommand\bkqjcup[5]{\begin{scope}[shift={#1}]
  \pgfmathsetmacro\bkqj@H{#3*(#4+0.6)}
  \fill[bkBlue!30] (-#2/2+0.6,0.6) rectangle (#2/2-0.6,{#3*#5});
  \draw[bkInk!80, line width=1pt, line join=round] (-#2/2,\bkqj@H) -- (-#2/2,0) -- (#2/2,0) -- (#2/2,\bkqj@H);
  \foreach \k in {1,...,#4} {\draw[bkInk, line width=0.6pt] (-#2/2,{#3*\k}) -- (-#2/2+3,{#3*\k});
    \node[anchor=east, font=\sffamily\bfseries\footnotesize, text=bkInk] at (-#2/2-0.8,{#3*\k}) {\k};}
\end{scope}}
% 目  \bkqjeye{(x,y)}{大きさ}
\newcommand\bkqjeye[2]{\begin{scope}[shift={#1}, scale=#2]
  \filldraw[fill=white, draw=bkInk, line width=0.7pt] (-1,0) .. controls (-0.4,0.62) and (0.4,0.62) .. (1,0) .. controls (0.4,-0.62) and (-0.4,-0.62) .. (-1,0) -- cycle;
  \fill[bkInk] (0.15,0) circle (0.33);\fill[white] (0.25,0.1) circle (0.09);
\end{scope}}
% 数字を大きくした はかり（第1巻の子ども向け）  \bkqjhakari{(x,y)}{半径}{1回りの量}{数字を書く間}{目もりの数}{針}{単位}
\newcommand\bkqjhakari[7]{\begin{scope}[shift={#1}]
  \fill[bkGray!35, rounded corners=2mm] (-#2-3,-#2-5) rectangle (#2+3,#2+3);
  \fill[bkGray!55] (-#2*0.8,#2+3) rectangle (#2*0.8,#2+5.5);
  \filldraw[fill=white, draw=bkInk!70, line width=0.8pt] (0,0) circle (#2);
  \foreach \k in {0,...,#5} {\pgfmathsetmacro\bkqj@v{#3/#5*\k}\pgfmathsetmacro\bkqj@a{90-\k/#5*360}
    \pgfmathtruncatemacro\bkqj@lab{mod(round(\bkqj@v),#4)==0}
    \ifnum\bkqj@lab=1 \draw[bkInk, line width=0.8pt] (\bkqj@a:#2*0.95) -- (\bkqj@a:#2*0.78);
    \else\draw[bkInk!80, line width=0.4pt] (\bkqj@a:#2*0.95) -- (\bkqj@a:#2*0.86);\fi}
  \pgfmathsetmacro\bkqj@fs{#2*0.5}%
  \foreach \k in {0,...,#5} {\pgfmathsetmacro\bkqj@v{#3/#5*\k}\pgfmathsetmacro\bkqj@a{90-\k/#5*360}
    \pgfmathtruncatemacro\bkqj@lab{and(mod(round(\bkqj@v),#4)==0, \k<#5)}\pgfmathtruncatemacro\bkqj@vi{round(\bkqj@v)}
    \ifnum\bkqj@lab=1 \node[font=\sffamily\bfseries\fontsize{\bkqj@fs}{\bkqj@fs}\selectfont, text=bkInk, inner sep=0pt] at (\bkqj@a:#2*0.58) {\bkqj@vi};\fi}
  \node[font=\sffamily\bfseries\fontsize{\bkqj@fs}{\bkqj@fs}\selectfont, text=bkGray] at (0,-#2*0.2) {#7};
  \pgfmathsetmacro\bkqj@na{90-#6/#3*360}
  \draw[bkRed, line width=1.3pt, line cap=round] (0,0) -- (\bkqj@na:#2*0.86);
  \fill[bkRed] (0,0) circle (1);
\end{scope}}
% 答えを 行の右に よせる  \bkqjright{中身}
\newcommand\bkqjright[1]{\unskip\nobreak\hfil\penalty50\hskip1em\null\nobreak\hfill\mbox{#1}\par}
\newcommand\bkqjans[1]{\bkqjright{{\sffamily\bfseries\small\color{bkOrange}こたえ}\enspace\bkkakikomi[16mm]\,#1}}
\makeatother

\bkchapter{1}{11}{\R{目}{め}もりを よもう}{\bklevel{1}}{\R{目}{め}もりを よもう}
\bkmichishirube{1-6,1-8,1-9}{1-16,1-17,2-2,2-6}

\begin{goal}
\begin{itemize}[label={\color{bkTheme}●}, leftmargin=1.4\zw, itemsep=0pt, topsep=0pt]
\item \mbox{1\R{目}{め}もり}が いくつかを \R{見}{み}つけられる。
\item \mbox{\R{温度計}{おん|ど|けい}}や はかりの \R{目}{め}もりを よめる。
\item \R{目}{め}もりの あいだを 「だいたい」で よめる。
\end{itemize}
\end{goal}

\begin{hatena}
\begin{bkfigure}
\begin{tikzpicture}[x=1mm, y=1mm]
  % A：温度計
  \fill[white, rounded corners=3mm] (0,0) rectangle (59,62);
  \node[circle, fill=bkYellow, text=white, font=\sffamily\bfseries, minimum size=6.5mm, inner sep=0pt] at (6,56) {A};
  \node[anchor=west, font=\small, text=bkInk] at (11,56) {\mbox{\R{温度計}{おん|ど|けい}}};
  \bkthermo{(38,12)}{0.95}{0}{40}{10}{26}
  \draw[bkOrange, line width=1pt, -{Stealth[length=4pt]}] (52,36.7) -- (43.5,36.7);
  \node[circle, draw=bkOrange, line width=1pt, fill=white, text=bkOrange!80!black, font=\sffamily\bfseries, inner sep=1pt, minimum size=6mm] at (53,40) {？};
  % B：はかり
  \begin{scope}[shift={(63,0)}]
  \fill[white, rounded corners=3mm] (0,0) rectangle (59,62);
  \node[circle, fill=bkYellow, text=white, font=\sffamily\bfseries, minimum size=6.5mm, inner sep=0pt] at (6,56) {B};
  \node[anchor=west, font=\small, text=bkInk] at (11,56) {はかり};
  \bkhakari{(28,24)}{18}{1000}{100}{50}{360}{g}
  \draw[bkOrange, line width=1pt, -{Stealth[length=4pt]}] (52,6) -- (41,12.5);
  \node[circle, draw=bkOrange, line width=1pt, fill=white, text=bkOrange!80!black, font=\sffamily\bfseries, inner sep=1pt, minimum size=6mm] at (53,6) {？};
  \end{scope}
\end{tikzpicture}
\end{bkfigure}
\noindent A の あかい ところは、 20と 30の あいだ。\\
B の はりは、 300と 400の あいだ。\\
\R{数字}{すう|じ}の ない \R{目}{め}もりは、 どう よむのかな？
\end{hatena}

\clearpage
%% ---------------- p.109 やってみよう ----------------
\begin{yattemiyou}[\R{目}{め}もりを かぞえよう]
\begin{enumerate}[label={\bkcircled[bkGreen]{\arabic*}}, leftmargin=2.2\zw, itemsep=2pt, topsep=0pt]
\item 10と 20の あいだは、 いくつに \R{分}{わ}かれて いるかな？ \R{目}{め}もりと \R{目}{め}もりの あいだを、 ゆびで おさえて かぞえよう。\enspace\bkkakikomi\enspace つ
\item 1つ すすむごとに 1ずつ ふやすと、 20の ところが 15に なって しまうね。
\item 2ずつ ふやすと どうかな？ \bkkakikomi[8mm]に かずを かこう。
\end{enumerate}
\begin{bkfigure}
\begin{tikzpicture}[x=1mm, y=1mm]
  % 目もり（10〜20、5つに分かれる）
  \fill[bkYellowBg, draw=bkYellow!80!black, line width=0.8pt, rounded corners=1mm] (-5,-10) rectangle (95,8);
  \foreach \k in {0,...,4} {\pgfmathtruncatemacro\bkqjo{mod(\k,2)}
    \ifnum\bkqjo=0 \fill[bkGreen!22] (\k*18+0.6,0.6) rectangle (\k*18+17.4,7.4);\fi}
  \draw[bkInk, line width=1.1pt] (0,8) -- (0,-3.5) (90,8) -- (90,-3.5);
  \foreach \k in {1,...,4} \draw[bkInk, line width=0.8pt] (\k*18,8) -- (\k*18,0);
  \node[font=\sffamily\bfseries\small] at (0,-6.8) {10};
  \node[font=\sffamily\bfseries\small] at (90,-6.8) {20};
  % 1ずつ
  \node[anchor=east, font=\sffamily\bfseries\footnotesize, text=bkGray] at (-6,-19) {1ずつ};
  \foreach \k/\v in {1/11,2/12,3/13,4/14} {\node[font=\sffamily\bfseries\small, text=bkGray] at (\k*18,-19) {\v};}
  \node[font=\sffamily\bfseries\small, text=bkRed] at (90,-19) {15};
  \node[font=\sffamily\bfseries\large, text=bkRed] at (101,-19) {×};
  % 2ずつ
  \node[anchor=east, font=\sffamily\bfseries\footnotesize, text=bkGreen!70!black] at (-6,-32) {2ずつ};
  \foreach \k/\v in {1/12,2/14} {\node[font=\sffamily\bfseries\small, text=bkGreen!60!black] at (\k*18,-32) {\v};}
  \foreach \k in {3,4} {\draw[bkTheme, line width=0.9pt, rounded corners=1mm, fill=white] (\k*18-5.5,-35.6) rectangle (\k*18+5.5,-28.4);}
  \node[font=\sffamily\bfseries\small, text=bkGreen!60!black] at (90,-32) {20};
  \draw[bkGreen!70!black, line width=1.4pt] (101,-32) circle (2.6);
  \foreach \y in {-12.5,-25.5} \draw[bkLine, line width=0.5pt, dash pattern=on 1.5pt off 1.5pt] (-16,\y) -- (105,\y);
\end{tikzpicture}
\end{bkfigure}
\vspace{4pt}
\noindent\hspace*{1\zw}2ずつ だと、 20に ぴったり！\\
\hspace*{1\zw}だから、 \R{目}{め}もり 1つ\R{分}{ぶん}は\enspace\bkkakikomi
\end{yattemiyou}

\bkfukidashi{mina}{\R{線}{せん}の かずじゃ なくて、\\ \R{線}{せん}と \R{線}{せん}の あいだを かぞえるんだね！}

\clearpage
%% ---------------- 数学の道具（第11回：1目もりの みつけかた・よみかたの 段を ふやした） ----------------
\bksection{doguu}{\mbox{\R{数学}{すう|がく}}の \mbox{\R{道具}{どう|ぐ}}}

\bksubsection{1}{\mbox{1\R{目}{め}もり}は いくつ？}

\label{bk:1-11:hitomemori}\R{目}{め}もり 1つ\R{分}{ぶん}の \R{大}{おお}きさを 「\mbox{1\R{目}{め}もり}」と いいます（1-6）。\bkidx{めもり}{\R{目}{め}もり} \R{目}{め}もりの あいだは どこも おなじ ながさなので、 1つ すすむごとに おなじ \R{数}{かず}ずつ ふえます。 \mbox{1\R{目}{め}もり}が いくつかは、 こう して \R{見}{み}つけます。
\begin{enumerate}[label={\bkcircled[bkTheme]{\arabic*}}, leftmargin=2.2\zw, itemsep=0pt, topsep=2pt]
\item \mbox{\R{数字}{すう|じ}}の ある \R{目}{め}もりを 2つ えらぶ。
\item 2つの \mbox{\R{数字}{すう|じ}}の ちがいを もとめる（1-5の ひき\R{算}{ざん}）。
\item その あいだが いくつに \R{分}{わ}かれて いるかを かぞえる。
\item 1ずつ、 2ずつ、 5ずつ、 10ずつ… と たして ためし、 ちがいに ぴったり なる \R{数}{かず}を さがす。
\end{enumerate}

\begin{bkfigure}
\begin{tikzpicture}[x=1mm, y=1mm]
  \bkqjscale{(5,0)}{10}{5}{5}{0}{100}
  \foreach \k in {1,...,5} {\node[circle, fill=bkGreen, text=white, font=\sffamily\bfseries\fontsize{7}{7}\selectfont, inner sep=0pt, minimum size=3.6mm] at (\k*10,7) {\k};}
  \node[anchor=west, font=\sffamily\bfseries\footnotesize, align=left] at (63,1) {ちがい 100 － 0 ＝ 100\\あいだは 5つ};
\end{tikzpicture}
\end{bkfigure}
\begin{center}
\begin{tblr}{colspec={Q[l,m]Q[l,m]Q[c,m]}, hlines={0.6pt,bkLine}, rowsep=2pt, colsep=6pt, row{1}={bg=bkGreen!12, font=\sffamily\bfseries\small}}
ためす \R{数}{かず} & 5かい たすと & 100に なる？ \\
10ずつ & 10、 20、 30、 40、 50 & \bkbatsu \\
20ずつ & 20、 40、 60、 80、 100 & \bkmaru \\
\end{tblr}
\end{center}
\noindent だから、 この \R{目}{め}もりの \mbox{1\R{目}{め}もり}は 20です。

\clearpage
\bksubsection{2}{\R{数字}{すう|じ}の ある \R{目}{め}もりから かぞえる}

\label{bk:1-11:kazoeru}\mbox{1\R{目}{め}もり}が わかったら、 はりや しるしの すぐ \R{左}{ひだり}（または \R{下}{した}）に ある \R{数字}{すう|じ}から、 その \R{数}{かず}ずつ かぞえます。

\begin{bkfigure}
\begin{tikzpicture}[x=1mm, y=1mm]
  \bkqjscale{(5,0)}{10}{10}{5}{0}{10}
  \bkqjmark{(85,3)}
  \foreach \k/\v in {1/12,2/14,3/16} {\draw[bkBlue, line width=1.1pt, -{Stealth[length=4pt]}] ({55+(\k-1)*10+0.6},-9.5) .. controls ({55+(\k-0.8)*10},-14.5) and ({55+(\k-0.2)*10},-14.5) .. ({55+\k*10-0.6},-9.7);
    \node[font=\sffamily\bfseries\scriptsize, text=bkBlue!85!black] at ({55+(\k-0.5)*10},-17) {\v};}
  \node[anchor=west, font=\sffamily\bfseries\footnotesize, text=bkRed!85!black] at (89,7) {ここは？};
\end{tikzpicture}
\end{bkfigure}
\begin{enumerate}[label={\bkcircled[bkTheme]{\arabic*}}, leftmargin=2.2\zw, itemsep=0pt, topsep=0pt]
\item 0と 10の あいだは 5つ。 2ずつで 10に ぴったり。 \mbox{1\R{目}{め}もり}は 2。
\item しるしの すぐ \R{左}{ひだり}の \R{数字}{すう|じ}は 10。
\item 10から 3つ すすむ。 12、 14、 16。 こたえは 16。
\end{enumerate}
\noindent 0から 2、 4、 6… と かぞえても 16に なりますが、 \R{近}{ちか}くの \R{数字}{すう|じ}から かぞえる ほうが はやくて、 まちがいも すくなく なります。

\bksubsection{3}{\mbox{\R{温度計}{おん|ど|けい}}を よもう}

\label{bk:1-11:ondokei}\mbox{\R{温度計}{おん|ど|けい}}は、 あたたかさを はかる \mbox{\R{道具}{どう|ぐ}}です。\bkidx{おんどけい}{\R{温度計}{おん|ど|けい}} \R{数直線}{すう|ちょく|せん}を たてに した ものと おなじで、 \R{上}{うえ}へ いくほど あたたかく なります。 あかい えきの \R{上}{うえ}の はしを よみます。

\noindent\begin{minipage}[c]{0.4\linewidth}\centering
\begin{tikzpicture}[x=1mm, y=1mm]
  \bkthermo{(14,0)}{1.4}{10}{30}{10}{26}
  \draw[bkTheme, line width=0.9pt, decorate, decoration={brace, amplitude=1.6mm}] (19.5,14) -- (19.5,22.4);
  \node[anchor=west, font=\sffamily\bfseries\footnotesize, text=bkTheme, align=left] at (22.5,18.2) {\mbox{6\R{目}{め}もり}\\\R{上}{うえ}};
\end{tikzpicture}
\end{minipage}\hfill
\begin{minipage}[c]{0.58\linewidth}
\begin{enumerate}[label={\bkcircled[bkTheme]{\arabic*}}, leftmargin=2\zw, itemsep=0pt, topsep=0pt]
\item 10と 20の あいだは 10に \R{分}{わ}かれて いる。 ちがいは 10なので、 \mbox{1\R{目}{め}もり}は 1。
\item えきの はしの すぐ \R{下}{した}の \R{数字}{すう|じ}は 20。
\item 20から \mbox{6\R{目}{め}もり} \R{上}{うえ}で、 26。
\end{enumerate}
\par\vspace{2pt}
こたえは 26\,℃（26ど）です。\bkidx{おんど}{\R{温度}{おん|ど}（℃）}
\end{minipage}
\par\vspace{4pt}
\noindent\bkkigocard{℃}{ど}{あたたかさや つめたさ （\mbox{\R{温度}{おん|ど}}） の たんい}{degree Celsius}{1-11}
\par\vspace{2pt}
\noindent\bkhatenaref の A は、 26\,℃です。

\bksubsection{4}{まるい \R{目}{め}もり}

\label{bk:1-11:marui}はかりの \R{目}{め}もりは、 まっすぐな \R{目}{め}もりを まるく した ものです。 0から 1まわりで、 もとの 0の ところに もどります。

\begin{bkfigure}
\begin{tikzpicture}[x=1mm, y=1mm]
  \bkqjscale{(4,2)}{3}{20}{5}{0}{100}
  \bkqjmark{(58,5)}
  \draw[bk arrow, bkGray, line width=1pt] (71,0) -- (83,0);
  \node[font=\sffamily\bfseries\footnotesize, text=bkGray] at (77,4) {まるく};
  \bkhakari{(108,0)}{18}{1000}{100}{50}{360}{g}
\end{tikzpicture}
\end{bkfigure}
\begin{enumerate}[label={\bkcircled[bkTheme]{\arabic*}}, leftmargin=2.2\zw, itemsep=0pt, topsep=0pt]
\item 300と 400の あいだを \R{見}{み}る。 ちがいは 100。
\item あいだは 5つ。 20ずつで 100に ぴったり。 \mbox{1\R{目}{め}もり}は 20\,g。
\item はりの すぐ \R{前}{まえ}の \R{数字}{すう|じ}は 300。 300から 320、 340、 360。
\end{enumerate}
\noindent\bkhatenaref の B は、 360\,g です。 まるい \R{目}{め}もりでは、 はりが まわる むき（\mbox{\R{時計}{と|けい}}の はりと おなじ むき）に すすむほど、 \R{数}{かず}が \R{大}{おお}きく なります。

\bksubsection{5}{あいだに ある ときは 「だいたい」}

\label{bk:1-11:daitai}はりや えきの はしが、 ちょうど \R{目}{め}もりに こない ときが あります。 \mbox{1\R{目}{め}もり}より こまかい ところは、 \R{目}{め}もりでは あらわせません。 そんな ときは、 \R{近}{ちか}い ほうの \R{目}{め}もりを よんで、 「だいたい」を つけます。\bkidx{だいたい}{だいたい}

\begin{bkfigure}
\begin{tikzpicture}[x=1mm, y=1mm]
  \fill[bkYellowBg, draw=bkYellow!80!black, line width=0.8pt, rounded corners=1mm] (-5,-9) rectangle (105,8);
  \fill[bkRed!85] (-5,1.8) rectangle (42.5,4.6);
  \foreach \k in {0,...,4} {\draw[bkInk, line width=0.8pt] (\k*25,8) -- (\k*25,-2.5);
    \pgfmathtruncatemacro\bkqjv{21+\k}\node[font=\sffamily\bfseries\small] at (\k*25,-5.8) {\bkqjv};}
  \draw[bkGray, line width=0.7pt, |-|] (25,12) -- (42.5,12);
  \node[font=\sffamily\bfseries\footnotesize, text=bkGray] at (33.75,15.5) {とおい};
  \draw[bkTheme, line width=0.9pt, |-|] (42.5,12) -- (50,12);
  \node[font=\sffamily\bfseries\footnotesize, text=bkTheme] at (50,15.5) {ちかい};
  \node[anchor=west, font=\sffamily\bfseries\small, text=bkTheme] at (66,13.5) {だいたい 23\,℃};
\end{tikzpicture}
\end{bkfigure}
\noindent \R{温度計}{おん|ど|けい}の えきの はしは、 22と 23の あいだ。 23の ほうが \R{近}{ちか}いので、 だいたい 23\,℃ です。


\bksubsection{6}{みの まわりの \R{目}{め}もり}

\label{bk:1-11:mawari}この \R{本}{ほん}で であった \R{目}{め}もりを まとめました。

\begin{center}
\begin{tblr}{colspec={Q[l,m,wd=40mm]Q[l,m,wd=30mm]Q[c,m,wd=34mm]}, hlines={0.6pt,bkLine}, vlines={0.6pt,bkLine}, rowsep=2pt, colsep=3pt,
  row{1}={bg=bkGreen!12, font=\sffamily\bfseries\small}}
\R{道具}{どう|ぐ} & はかる もの & \mbox{1\R{目}{め}もり} \\
ものさし（1-8） & ながさ & 1\,mm \\
1\,L ます（1-9） & かさ & 1\,dL \\
\mbox{\R{時計}{と|けい}}（1-10） & \R{時}{じ}こく・\R{時間}{じ|かん} & \mbox{1\R{分}{ぷん}} \\
\R{温度計}{おん|ど|けい}（この \R{章}{しょう}） & あたたかさ & 1\,℃ \\
はかり（この \R{章}{しょう}） & おもさ & 20\,g など \\
\end{tblr}
\end{center}

\noindent{\sffamily\bfseries\color{bkTheme}さがして みよう}\quad いえの \R{中}{なか}で \R{目}{め}もりの ある ものを さがして、 \mbox{1\R{目}{め}もり}を しらべよう。 おなじ しゅるいの \R{道具}{どう|ぐ}でも、 \mbox{1\R{目}{め}もり}は ちがう ことが あります。

\begin{taisetsu}
\mbox{1\R{目}{め}もり}が いくつかを \R{見}{み}つけてから、 \R{近}{ちか}くの \R{数字}{すう|じ}から かぞえる。\\
\R{目}{め}もりの あいだは、 \R{近}{ちか}い ほうを よんで 「だいたい」。
\end{taisetsu}

\clearpage
%% ---------------- 例題 ----------------
\begin{reidai}[\mbox{1\R{目}{め}もり}は いくつ？]
この \R{目}{め}もりの \mbox{1\R{目}{め}もり}は いくつ？
\begin{bkfigure}
\begin{tikzpicture}[x=1mm, y=1mm]
  \bkqjscale{(0,0)}{9}{10}{5}{0}{50}
\end{tikzpicture}
\end{bkfigure}
\kangaekata 0と 50の ちがいは 50。 あいだは 5つ。 5ずつだと 5、 10、 15、 20、 25で ×。 10ずつだと 10、 20、 30、 40、 50で ぴったり。 \kotaenext 10
\end{reidai}

\begin{jibun}
\noindent\begin{minipage}[c]{0.6\linewidth}\centering
\begin{tikzpicture}[x=1mm, y=1mm]
  \bkqjscale{(0,0)}{7}{8}{4}{0}{20}
\end{tikzpicture}
\end{minipage}\hfill
\begin{minipage}[c]{0.38\linewidth}\raggedleft
\mbox{1\R{目}{め}もり}は いくつ？\par\vspace{2pt}
{\sffamily\bfseries\small\color{bkOrange}こたえ}\enspace\bkkakikomi[16mm]
\end{minipage}
\bkjans{5}
\end{jibun}

\begin{reidai}[\mbox{1\R{目}{め}もり}が 2の \R{目}{め}もりを よむ]
\noindent\begin{minipage}[c]{0.4\linewidth}
しるしの ところは いくつ？
\end{minipage}\hfill
\begin{minipage}[c]{0.58\linewidth}\raggedleft
\begin{tikzpicture}[x=1mm, y=1mm]
  \bkqjscale{(0,0)}{6}{10}{5}{30}{10}
  \bkqjmark{(42,3)}
\end{tikzpicture}
\end{minipage}
\par\vspace{2pt}
\kangaekata 30と 40の ちがいは 10、 あいだは 5つ。 2ずつで 10に ぴったりだから、 \mbox{1\R{目}{め}もり}は 2。 しるしの \R{左}{ひだり}の \R{数字}{すう|じ}は 40。 40から 2つで、 42、 44。 \kotaenext 44
\end{reidai}

\begin{jibun}
\noindent\begin{minipage}[c]{0.6\linewidth}\centering
\begin{tikzpicture}[x=1mm, y=1mm]
  \bkqjscale{(0,0)}{6}{10}{5}{60}{10}
  \bkqjmark{(18,3)}
\end{tikzpicture}
\end{minipage}\hfill
\begin{minipage}[c]{0.38\linewidth}\raggedleft
しるしの ところは いくつ？\par\vspace{2pt}
{\sffamily\bfseries\small\color{bkOrange}こたえ}\enspace\bkkakikomi[16mm]
\end{minipage}
\bkjans{66（60から 62、 64、 66）}
\end{jibun}

\begin{reidai}[\mbox{\R{温度計}{おん|ど|けい}}を よむ]
\noindent\begin{minipage}[c]{0.76\linewidth}
この \mbox{\R{温度計}{おん|ど|けい}}は なん\,℃？

\kangaekata 10と 20の あいだは 10に \R{分}{わ}かれて いるので、 \mbox{1\R{目}{め}もり}は 1\,℃。\\
えきの はしの \R{下}{した}の \R{数字}{すう|じ}は 10。 10から \mbox{7\R{目}{め}もり} \R{上}{うえ}。

\kotae 17\,℃
\end{minipage}\hfill
\begin{minipage}[c]{0.22\linewidth}\centering
\begin{tikzpicture}[x=1mm, y=1mm]
  \bkthermo{(8,0)}{1.2}{10}{30}{10}{17}
\end{tikzpicture}
\end{minipage}
\end{reidai}

\begin{jibun}
\noindent\begin{minipage}[c]{0.76\linewidth}
この \mbox{\R{温度計}{おん|ど|けい}}は なん\,℃？\par
{\small\color{bkOrange!58!black}ヒント：\R{数字}{すう|じ}は 5ごと。 20と 25の あいだは いくつに \R{分}{わ}かれて いる？}\bkqjans{℃}
\end{minipage}\hfill
\begin{minipage}[c]{0.22\linewidth}\centering
\begin{tikzpicture}[x=1mm, y=1mm]
  \bkthermo{(8,0)}{1.3}{20}{30}{5}{24}
\end{tikzpicture}
\end{minipage}
\bkjans{24\,℃（\mbox{1\R{目}{め}もり}は 1\,℃）}
\end{jibun}

\begin{reidai}[はかりを よむ]
\noindent\begin{minipage}[c]{0.64\linewidth}
この はかりは なん\,g？

\kangaekata 0と 100の あいだは 5つ。 20ずつで 100に ぴったりなので、 \mbox{1\R{目}{め}もり}は 20\,g。\\
はりの \R{前}{まえ}の \R{数字}{すう|じ}は 300。 300から 320、 340。

\kotae 340\,g
\end{minipage}\hfill
\begin{minipage}[c]{0.34\linewidth}\centering
\begin{tikzpicture}[x=1mm, y=1mm]
  \bkqjhakari{(0,0)}{12}{500}{100}{25}{340}{g}
\end{tikzpicture}
\end{minipage}
\par\vspace{2pt}
\noindent{\small まっすぐに のばして \R{見}{み}ると：}
\begin{bkfigure}
\begin{tikzpicture}[x=1mm, y=1mm]
  \bkqjscale{(5,0)}{14}{5}{5}{300}{100}
  \foreach \k/\v in {1/320,2/340,3/360,4/380} {\node[font=\sffamily\bfseries\scriptsize, text=bkBlue!85!black] at (5+\k*14,-5) {\v};}
  \bkqjmark{(33,3)}
\end{tikzpicture}
\end{bkfigure}
\end{reidai}

\begin{jibun}
\noindent\begin{minipage}[c]{0.64\linewidth}
この はかりは なん\,g？\par
{\small\color{bkOrange!58!black}ヒント：0と 100の あいだは いくつに \R{分}{わ}かれて いる？}\par\vspace{4pt}
\noindent\hfill{\sffamily\bfseries\small\color{bkOrange}こたえ}\enspace\bkkakikomi[16mm]\,g
\end{minipage}\hfill
\begin{minipage}[c]{0.34\linewidth}\centering
\begin{tikzpicture}[x=1mm, y=1mm]
  \bkqjhakari{(0,0)}{12}{500}{100}{25}{180}{g}
\end{tikzpicture}
\end{minipage}
\bkjans{180\,g（100から 120、 140、 160、 180）}
\end{jibun}

\begin{reidai}[\mbox{1\R{目}{め}もり}が 5の メーター]
\noindent\begin{minipage}[c]{0.6\linewidth}
はりは いくつを さして いる？

\kangaekata 0と 20の あいだは 4つ。 5ずつで 5、 10、 15、 20と ぴったり。 \mbox{1\R{目}{め}もり}は 5。\\
はりの \R{前}{まえ}の \R{数字}{すう|じ}は 40。 40から 1つで 45。

\kotae 45
\end{minipage}\hfill
\begin{minipage}[c]{0.38\linewidth}\centering
\begin{tikzpicture}[x=1mm, y=1mm]
  \bkqjgauge{(0,0)}{21}{100}{20}{20}{45}
\end{tikzpicture}
\end{minipage}
\end{reidai}

\begin{jibun}
\noindent\begin{minipage}[c]{0.6\linewidth}
はりは いくつを さして いる？\par\vspace{8pt}
\noindent\hfill{\sffamily\bfseries\small\color{bkOrange}こたえ}\enspace\bkkakikomi[16mm]
\end{minipage}\hfill
\begin{minipage}[c]{0.38\linewidth}\centering
\begin{tikzpicture}[x=1mm, y=1mm]
  \bkqjgauge{(0,0)}{20}{100}{20}{20}{85}
\end{tikzpicture}
\end{minipage}
\bkjans{85（80から 1つ）}
\end{jibun}

\begin{reidai}[だいたい いくつ？]
\noindent\begin{minipage}[c]{0.7\linewidth}
ジュースは だいたい なん\,dL？

\kangaekata \R{水}{みず}の \R{上}{うえ}の \R{面}{めん}は、 3と 4の あいだ。 4の ほうに \R{近}{ちか}い。

\kotae だいたい 4\,dL
\end{minipage}\hfill
\begin{minipage}[c]{0.28\linewidth}\centering
\begin{tikzpicture}[x=1mm, y=1mm]
  \bkqjcup{(0,0)}{18}{4.8}{5}{3.7}
  \node[font=\sffamily\bfseries\footnotesize, text=bkGray] at (0,30) {dL};
\end{tikzpicture}
\end{minipage}
\end{reidai}

\begin{jibun}
\noindent\begin{minipage}[c]{0.46\linewidth}\centering
\begin{tikzpicture}[x=1mm, y=1mm]
  \bkqjscale{(0,0)}{4.6}{10}{1}{0}{1}
  \bkpencil{(2.2,7)}{29.5}{bkRoute4}
\end{tikzpicture}
\end{minipage}\hfill
\begin{minipage}[c]{0.52\linewidth}\raggedleft
だいたい なん\,cm？\par\vspace{2pt}
{\sffamily\bfseries\small\color{bkOrange}こたえ}\enspace\bkkakikomi[14mm]\,cm
\end{minipage}
\bkjans{だいたい 7\,cm}
\end{jibun}

\begin{reidai}[\R{目}{め}もりを よんで くらべる]
\noindent\begin{minipage}[c]{0.68\linewidth}
あさと ひるの \R{気温}{き|おん}です。 なん\,℃ あがった？

\kangaekata \mbox{1\R{目}{め}もり}は 1\,℃。 あさは 10から 8つで 18\,℃、 ひるは 20から 5つで 25\,℃。 あがった \R{分}{ぶん}は ちがいなので ひき\R{算}{ざん}（1-5、 1-7）。

\shiki 25 － 18 ＝ 7 \kotaenext 7\,℃
\end{minipage}\hfill
\begin{minipage}[c]{0.3\linewidth}\centering
\begin{tikzpicture}[x=1mm, y=1mm]
  \bkthermo{(9,0)}{1.1}{10}{30}{10}{18}
  \bkthermo{(30,0)}{1.1}{10}{30}{10}{25}
  \node[anchor=north, font=\sffamily\bfseries\scriptsize] at (9,-9.5) {あさ};
  \node[anchor=north, font=\sffamily\bfseries\scriptsize] at (30,-9.5) {ひる};
\end{tikzpicture}
\end{minipage}
\end{reidai}

\begin{jibun}
\noindent\begin{minipage}[c]{0.68\linewidth}
\R{水}{みず}を あたためました。 なん\,℃ あがった？
\bkshikiran[6\zw]
\end{minipage}\hfill
\begin{minipage}[c]{0.3\linewidth}\centering
\begin{tikzpicture}[x=1mm, y=1mm]
  \bkthermo{(9,0)}{1.1}{10}{30}{10}{12}
  \bkthermo{(30,0)}{1.1}{10}{30}{10}{21}
  \node[anchor=north, font=\sffamily\bfseries\scriptsize] at (9,-9.5) {まえ};
  \node[anchor=north, font=\sffamily\bfseries\scriptsize] at (30,-9.5) {あと};
\end{tikzpicture}
\end{minipage}
\bkjans{21 － 12 ＝ 9\quad 9\,℃（まえ 12\,℃、 あと 21\,℃）}
\end{jibun}


\clearpage
%% ---------------- p.114 まちがえやすいところ ----------------
\bksection{tsumazuki}{まちがえやすい ところ}

\begin{tsumazuki}{1ずつ よんで しまう}
\noindent\begin{minipage}[c]{0.6\linewidth}
{\color{bkRed}▼}は？\hspace{1\zw}\bkbatsu\enspace 303\,g\hspace{2\zw}\bkmaru\enspace 360\,g\par
\mbox{1\R{目}{め}もり}は 20\,g。 \mbox{1\R{目}{め}もり}が いくつかを、 さきに \R{見}{み}つけよう。
\end{minipage}\hfill
\begin{minipage}[c]{0.38\linewidth}\centering
\begin{tikzpicture}[x=1mm, y=1mm]
  \bkqjscale{(0,0)}{7}{5}{5}{300}{100}
  \bkqjmark{(21,3)}
\end{tikzpicture}
\end{minipage}
\end{tsumazuki}

\begin{tsumazuki}{\R{線}{せん}を かぞえて しまう}
\noindent\begin{minipage}[c]{0.6\linewidth}
10と 20の あいだの \R{線}{せん}は \mbox{4\R{本}{ほん}}。 でも、 あいだは 5つです。 あいだを かぞえよう。
\end{minipage}\hfill
\begin{minipage}[c]{0.38\linewidth}\centering
\begin{tikzpicture}[x=1mm, y=1mm]
  \bkqjscale{(0,0)}{7}{5}{5}{10}{10}
  \foreach \k in {0,...,4} {\node[circle, fill=bkGreen, text=white, font=\sffamily\bfseries\fontsize{7}{7}\selectfont, inner sep=0pt, minimum size=3.6mm] at (\k*7+3.5,7) {\the\numexpr\k+1\relax};}
\end{tikzpicture}
\end{minipage}
\end{tsumazuki}

\begin{tsumazuki}{\R{目}{め}の \R{高}{たか}さ}
\noindent\begin{minipage}[c]{0.6\linewidth}
\R{上}{うえ}や \R{下}{した}から \R{見}{み}ると、 ずれて よんで しまいます。 \R{目}{め}の \R{高}{たか}さを、 \R{水}{みず}の \R{高}{たか}さに あわせよう。
\end{minipage}\hfill
\begin{minipage}[c]{0.38\linewidth}\centering
\begin{tikzpicture}[x=1mm, y=1mm]
  \bkqjcup{(10,0)}{14}{6}{4}{2}
  \bkqjeye{(34,24)}{2.6}\bkqjeye{(34,12)}{2.6}\bkqjeye{(34,0)}{2.6}
  \draw[bkRed!70, line width=0.6pt, dash pattern=on 1.5pt off 1pt] (30.5,24) -- (17,12);
  \draw[bkGreen!70!black, line width=0.8pt, dash pattern=on 1.5pt off 1pt] (30.5,12) -- (17,12);
  \draw[bkRed!70, line width=0.6pt, dash pattern=on 1.5pt off 1pt] (30.5,0) -- (17,12);
  \node[font=\sffamily\bfseries\small, text=bkRed] at (41,24) {×};
  \node[font=\sffamily\bfseries\small, text=bkGreen!70!black] at (41,12) {○};
  \node[font=\sffamily\bfseries\small, text=bkRed] at (41,0) {×};
\end{tikzpicture}
\end{minipage}
\end{tsumazuki}

\begin{column}{℃の はじまり}
\R{水}{みず}が こおる ときを 0\,℃、 \R{水}{みず}が にえたつ ときを 100\,℃に して、 その あいだを 100に \R{分}{わ}けた \R{目}{め}もりです。 セルシウスと いう \R{人}{ひと}の \R{名前}{な|まえ}から、 ℃と かきます。
\end{column}

\clearpage
%% ---------------- p.115 つかわれているところ ----------------
\begin{butsuri}
\R{目}{め}もりを よむ ことは、 きかいを まもる しごとでも だいじです。
\begin{bkfigure}
\begin{tikzpicture}[x=1mm, y=1mm]
  % 工場のメーター
  \fill[white, rounded corners=3mm] (0,0) rectangle (59,40);
  \fill[bkGray!55] (3,11) rectangle (34,16);
  \fill[bkGray!70] (13,16) rectangle (15,22);
  \fill[bkGray!35, rounded corners=1.5mm] (3,21) rectangle (25,36);
  \filldraw[fill=white, draw=bkInk!70, line width=0.7pt] (5,24) arc (180:0:9) -- cycle;
  \foreach \k in {0,...,10} {\draw[bkInk, line width=0.4pt] ($(14,24)+({180-\k*18}:8.6)$) -- ($(14,24)+({180-\k*18}:7.2)$);}
  \draw[bkGreen, line width=1.6pt] (14,24) ++(180:9.6) arc (180:72:9.6);
  \draw[bkRed, line width=1.6pt] (14,24) ++(36:9.6) arc (36:0:9.6);
  \draw[bkRed, line width=1pt, line cap=round] (14,24) -- ++(108:7.4);
  \fill[bkInk!80] (14,24) circle (1);
  % ロボットのカメラ
  \fill[bkNavy!75, rounded corners=1mm] (38,24) rectangle (52,32);
  \fill[white] (41.5,28) circle (2.4); \fill[bkInk] (41.5,28) circle (1.3);
  \fill[bkNavy!60] (43.5,9) rectangle (46.5,24);
  \fill[bkNavy!75] (39,7) rectangle (51,9);
  \draw[bkOrange, line width=0.7pt, dash pattern=on 1.5pt off 1pt] (38,29.5) -- (24,32) (38,26.5) -- (24,25);
  \node[font=\sffamily\bfseries\footnotesize] at (29.5,3.5) {\R{工場}{こう|じょう}の メーター};
  % 車のメーター
  \begin{scope}[shift={(63,0)}]
  \fill[white, rounded corners=3mm] (0,0) rectangle (59,40);
  \fill[bkInk!85, rounded corners=3mm] (3,10) rectangle (56,36);
  \begin{scope}[shift={(20,17)}]
    \fill[white] (-12,0) arc (180:0:12) -- cycle;
    \foreach \k in {0,...,12} {\draw[bkInk, line width=0.5pt] ({180-\k*15}:11.6) -- ({180-\k*15}:10);}
    \foreach \k/\v in {3/30,6/60,9/90} {\node[font=\sffamily\bfseries\fontsize{6.5}{7}\selectfont, inner sep=0pt] at ({180-\k*15}:7.6) {\v};}
    \draw[bkRed, line width=1.1pt, line cap=round] (0,0) -- (120:10);
    \fill[bkInk] (0,0) circle (1);
  \end{scope}
  \begin{scope}[shift={(45,17)}]
    \fill[white] (-8,0) arc (180:0:8) -- cycle;
    \node[font=\sffamily\bfseries\fontsize{7}{7}\selectfont, text=bkRed] at (-5.5,2) {E};
    \node[font=\sffamily\bfseries\fontsize{7}{7}\selectfont, text=bkInk] at (5.5,2) {F};
    \draw[bkRed, line width=1.1pt, line cap=round] (0,0) -- (60:7);
    \fill[bkInk] (0,0) circle (0.9);
  \end{scope}
  \node[font=\sffamily\bfseries\fontsize{8}{9}\selectfont, text=white] at (20,12.5) {はやさ};
  \node[font=\sffamily\bfseries\fontsize{8}{9}\selectfont, text=white] at (45,12.5) {ガソリン};
  \node[font=\sffamily\bfseries\footnotesize] at (29.5,3.5) {\R{車}{くるま}の メーター};
  \end{scope}
\end{tikzpicture}
\end{bkfigure}
\begin{itemize}[label={\color{bkPurple}●}, leftmargin=1.4\zw, itemsep=0pt, topsep=0pt]
\item \R{工場}{こう|じょう}の パイプには、 まるい メーターが たくさん あります。 \R{人}{ひと}や ロボットが はりを よんで、 いつもと ちがう ときは すぐに しらせます。
\item \R{車}{くるま}の メーターは、 はやさや ガソリンの のこりを はりで しめします。
\item \R{天気}{てん|き}よほうの 「きょうは 25\,℃」や、 \R{家}{いえ}の \R{電気}{でん|き}の メーターも、 はかった かずです。
\end{itemize}
\end{butsuri}

\bkfukidashi{delta}{はかった かずを、 いつもの かずと くらべると、\\ きかいの 「びょうき」に はやく \R{気}{き}づけるよ。}

\begin{bkfigure}
\begin{tikzpicture}[x=1mm, y=1mm]
  \draw[bkRoute7, line width=3pt] (10,0) -- (128,0);
  \node[anchor=west, inner sep=0pt] at (0,12) {\bkrouteicon[1.4em]{7}\hspace{0.3em}{\sffamily\bfseries\small\color{bkRoute7!80!black}「はかる」の \R{路線}{ろ|せん}}};
  \foreach \x/\c/\t/\s in {10/bkVol1/{1-8}/{ながさ},33.6/bkVol1/{1-9}/{かさ・おもさ},57.2/bkVol1/{1-11}/{\R{目}{め}もり},80.8/bkVol1/{1-16}/{\R{小数}{しょう|すう}},104.4/bkVol2/{2-6}/{ずれ},128/bkVol3/{3-14}/{ばらつき}} {
    \fill[white, draw=\c, line width=1.6pt] (\x,0) circle (2.6);
    \node[anchor=base, font=\sffamily\bfseries\footnotesize, text=\c!85!black] at (\x,-8.4) {\t};
    \node[anchor=base, font=\sffamily\footnotesize, text=bkInk] at (\x,-15.4) {\s};}
  \node[fill=bkTheme, text=white, rounded corners=1mm, font=\sffamily\bfseries\fontsize{8}{9}\selectfont, inner sep=1.5pt] at (57.2,5.5) {いま ここ};
\end{tikzpicture}
\end{bkfigure}

\clearpage
%% ---------------- p.116〜117 もんだい ----------------
\bkendproblems

\bkproblemgroup{きほん}{1}

\bkproblem \mbox{1\R{目}{め}もり}は いくつ？
\begin{bkfigure}
\begin{tikzpicture}[x=1mm, y=1mm]
  \foreach \i/\d/\n/\st in {0/3.2/10/10,1/6.4/5/10,2/8/4/20} {
    \node[font=\sffamily\bfseries\small] at (\i*45-1,8) {（\the\numexpr\i+1\relax）};
    \bkqjscale{(\i*45+4,0)}{\d}{\n}{\n}{0}{\st}
    \draw[bkTheme, line width=0.9pt, rounded corners=1mm, fill=white] (\i*45+13,-17) rectangle ++(14,7.2);}
\end{tikzpicture}
\end{bkfigure}
\bkpans{\mbox{（1）1}\quad \mbox{（2）2}\quad \mbox{（3）5}}

\bkproblem \mbox{\R{温度計}{おん|ど|けい}}は なん\,℃？
\begin{bkfigure}
\begin{tikzpicture}[x=1mm, y=1mm]
  \foreach \i/\t in {0/15,1/28} {
    \node[font=\sffamily\bfseries\small] at (\i*62,30) {（\the\numexpr\i+1\relax）};
    \bkthermo{(\i*62+20,4)}{1.1}{10}{30}{10}{\t}
    \draw[bkTheme, line width=0.9pt, rounded corners=1mm, fill=white] (\i*62+30,4) rectangle ++(16,7.2);
    \node[anchor=west, font=\small] at (\i*62+46.5,7.6) {℃};}
\end{tikzpicture}
\end{bkfigure}
\bkpans{\mbox{（1）15\,℃}\quad \mbox{（2）28\,℃}}

\noindent\begin{minipage}[t]{0.49\linewidth}
\bkproblem はかりは なん\,g？\par\vspace{2pt}
\noindent\begin{minipage}[c]{0.56\linewidth}\centering
\begin{tikzpicture}[x=1mm, y=1mm]
  \bkqjhakari{(0,0)}{13}{500}{100}{25}{460}{g}
\end{tikzpicture}
\end{minipage}\hfill
\begin{minipage}[c]{0.42\linewidth}
\bkkakikomi[14mm]\,g
\end{minipage}
\bkpans{460\,g}
\end{minipage}\hfill
\begin{minipage}[t]{0.49\linewidth}
\bkproblem \R{水}{みず}は だいたい なん\,dL？\par\vspace{2pt}
\noindent\begin{minipage}[c]{0.4\linewidth}\centering
\begin{tikzpicture}[x=1mm, y=1mm]
  \bkqjcup{(0,0)}{16}{5.5}{5}{2.2}
  \node[font=\sffamily\bfseries\footnotesize, text=bkGray] at (0,33) {dL};
\end{tikzpicture}
\end{minipage}\hfill
\begin{minipage}[c]{0.58\linewidth}
だいたい\\\bkkakikomi[14mm]\,dL
\end{minipage}
\bkpans{だいたい 2\,dL}
\end{minipage}

\clearpage
\bkproblemgroup{チャレンジ}{1}

\bkproblem メーターを みましょう。
\par\vspace{2pt}
\noindent\begin{minipage}[c]{0.6\linewidth}
（1）\mbox{1\R{目}{め}もり}は いくつ？\par
{\small\color{bkOrange!58!black}ヒント：20と 40の あいだを かぞえよう}\par\vspace{2pt}
\noindent\hfill\bkkakikomi[16mm]\par\vspace{4pt}
（2）はりは いくつを さして いる？\par\vspace{2pt}
\noindent\hfill\bkkakikomi[16mm]
\end{minipage}\hfill
\begin{minipage}[c]{0.38\linewidth}\centering
\begin{tikzpicture}[x=1mm, y=1mm]
  \bkqjgauge{(0,0)}{21}{100}{20}{20}{65}
\end{tikzpicture}
\end{minipage}
\bkpans{\mbox{（1）5}\quad \mbox{（2）65}}

\bkproblem \R{体重計}{たい|じゅう|けい}は なん\,kg？
\begin{bkfigure}
\begin{tikzpicture}[x=1mm, y=1mm]
  \fill[bkGray!30, rounded corners=3mm] (-10,-12) rectangle (100,12);
  \bkqjscale{(0,0)}{6}{15}{5}{20}{10}
  \bkqjmark{(42,3)}
  \draw[bkTheme, line width=0.9pt, rounded corners=1mm, fill=white] (104,-3.6) rectangle ++(16,7.2);
  \node[anchor=west, font=\small] at (120.5,0) {kg};
\end{tikzpicture}
\end{bkfigure}
\bkpans{34\,kg}

\begin{furikaeri}
\begin{checklist}
\item \mbox{1\R{目}{め}もり}が いくつかを \R{見}{み}つけられる。\hfill{\footnotesize\sffamily\color{bkTheme}→\pageref{bk:1-11:hitomemori}ページ}
\item \mbox{\R{温度計}{おん|ど|けい}}や はかりの \R{目}{め}もりを よめる。\hfill{\footnotesize\sffamily\color{bkTheme}→\pageref{bk:1-11:ondokei}ページ}
\item \R{目}{め}もりの あいだを 「だいたい」で よめる。\hfill{\footnotesize\sffamily\color{bkTheme}→\pageref{bk:1-11:daitai}ページ}
\end{checklist}
チェックが つかなかったら、 →の ページに もどろう。
\end{furikaeri}

\bkfukidashi{sora}{これで 「はかる」の \R{部}{ぶ}は おしまい。\\ つぎは、 かけ\R{算}{ざん}の \R{部}{ぶ}だよ！}
\clearpage
